\documentclass[12pt]{report}
\usepackage{setspace}
\usepackage{minitoc}
\usepackage{amsmath,amssymb}
\newcommand{\thesisspacing}{\doublespacing}
\dominitoc

\begin{document}

\doublespacing % Do not change - required

\chapter{Implementation Details}
\label{ch1}

%%%%%%%%%%%%%%%%%%%%%%%%%%%%%%%%%%%%%%%
% IMPORTANT
\begin{spacing}{1} %THESE FOUR
\minitoc % LINES MUST APPEAR IN
\end{spacing} % EVERY
\thesisspacing % CHAPTER
% COPY THEM IN ANY NEW CHAPTER

\section{Overview}

This chapter describes the implementation of HRAC (History-Residual Adaptive Clipping). HRAC is designed to improve Byzantine robustness while preserving model utility in non-IID federated settings.

\section{Core Design}

HRAC consists of two core mechanisms:
\begin{enumerate}
    \item \textbf{Global Norm Clipping}
    \item \textbf{$\nu$-based Soft Downweighting}
\end{enumerate}

\section{Pipeline}

At round $t$, with client updates $\{\Delta_i\}_{i=1}^{N}$:
\begin{enumerate}
    \item Compute per-client norms and their median.
    \item Apply global clipping to obtain $\Delta_i^{\mathrm{clip}}$.
    \item Build residuals using client history and construct $\tilde{\Delta}_i$.
    \item After warm-up, compute normalized $\nu$-based weights $w_i$.
    \item Aggregate: $g_t=\sum_i w_i\tilde{\Delta}_i$.
    \item Update client history states.
\end{enumerate}

\section{Global Norm Clipping}

\begin{equation}
\mathrm{med}_t = \mathrm{median}\left(\{\|\Delta_i\|\}_{i=1}^{N}\right), \qquad
B_t = c_g \cdot \mathrm{med}_t.
\end{equation}

\begin{equation}
s_i^{\mathrm{global}}=\min\left(1,\frac{B_t}{\|\Delta_i\|+\epsilon}\right), \qquad
\Delta_i^{\mathrm{clip}}=s_i^{\mathrm{global}}\Delta_i.
\end{equation}

\section{History-Residual Modeling}

\begin{equation}
r_i=\Delta_i^{\mathrm{clip}}-b_i, \qquad \tau_i=c\cdot\mu_i+\epsilon.
\end{equation}

In the current HRAC run, per-client residual clipping is always enabled:
\begin{equation}
\bar{r}_i=\mathrm{clip\_by\_norm}(r_i,\tau_i),
\end{equation}
and the candidate update is
\begin{equation}
\tilde{\Delta}_i=b_i+\bar{r}_i.
\end{equation}

Following the current experiment configuration (\texttt{main CMomentum.py}), no extra post-reconstruction cap is applied to $\tilde{\Delta}_i$.

\subsection*{How $b_i$ is obtained}

The history mean $b_i$ is initialized and updated as follows.

For a new client $i$ (first appearance in training),
\begin{equation}
b_i \leftarrow \Delta_i^{\mathrm{clip}}.
\end{equation}

For subsequent rounds, after aggregation-path processing, $b_i$ is updated by EMA:
\begin{equation}
\Delta_i^{\mathrm{safe}}=\tilde{\Delta}_i\cdot
\min\left(1,\frac{B_t}{\|\tilde{\Delta}_i\|+\epsilon}\right),
\end{equation}
\begin{equation}
b_i \leftarrow \rho_b\, b_i + (1-\rho_b)\,\Delta_i^{\mathrm{safe}}.
\end{equation}

Therefore, $b_i$ is not a fixed baseline; it is a client-specific moving center that tracks stable update behavior over rounds.

\section{$\nu$ Statistics and Downweighting}

\begin{equation}
d_i=\left\|\bar{r}_i^{\mathrm{norm}}-r_{i,\mathrm{prev}}\right\|, \qquad
\nu_i \leftarrow \rho_\nu \nu_i + (1-\rho_\nu)d_i.
\end{equation}

After warm-up:
\begin{equation}
e_i=
\begin{cases}
\max(\nu_i-\bar{\nu},0), & \nu_i>1,\\
0, & \nu_i \le 1,
\end{cases}
\end{equation}
\begin{equation}
\hat{w}_i=\frac{1}{1+\alpha e_i}, \qquad
w_i=\frac{\hat{w}_i}{\sum_j \hat{w}_j+\epsilon}.
\end{equation}

The per-client cap \texttt{nu\_weight\_max} is enforced by clipping and renormalization.

\section{Aggregation and State Updates}

\begin{equation}
g_t=\sum_{i=1}^{N}w_i\tilde{\Delta}_i.
\end{equation}

Main maintained states are \texttt{b[cid]}, \texttt{b\_norm[cid]}, \texttt{mu[cid]}, \texttt{nu[cid]}, \texttt{r\_prev[cid]}, and \texttt{median\_norm\_prev}.

\section{Complexity}

With $n$ clients and model dimension $d$:
\begin{equation}
\text{Time}=\mathcal{O}(nd), \qquad \text{Memory}=\mathcal{O}(nd).
\end{equation}

\section{Summary}

HRAC combines global norm clipping, history-residual modeling, and $\nu$-based soft downweighting in a unified aggregation pipeline for robust non-IID federated learning.

\end{document}
